% ============================================================
% OBJECTIVES & SOLUTION SPECIFICATION
% ============================================================

\section{Engineering Design Objectives \& Solution Specifications}

\subsection{Primary Functional Goals}

List the primary engineering objectives.

\begin{itemize}
    \item \textbf{Objective 1:} Describe the first measurable objective.
    \item \textbf{Objective 2:} Describe the second measurable objective.
    \item \textbf{Objective 3:} Describe the third measurable objective.
\end{itemize}

\subsection{High-Level Architectural Workflow}

Describe the interaction between hardware, software, communication,
processing, and user interfaces.

\begin{figure}[h!]
    \centering
    \begin{tikzpicture}[node distance=1.5cm, auto]

        \node [hlddstep] (s1) {Input / Data Generation};
        \node [hlddstep] (s2) [below=0.8cm of s1] {Processing / Communication};
        \node [hlddstep] (s3) [below=0.8cm of s2] {Analysis / Decision};
        \node [hldddecision] (d1) [below=0.8cm of s3] {Condition Met?};
        \node [hlddalert] (a1) [right=1.5cm of d1] {Alert / Corrective Action};
        \node [hlddstep] (s4) [below=1.2cm of d1] {Normal Operation / Output};

        \draw [hlddarrow] (s1) -- (s2);
        \draw [hlddarrow] (s2) -- (s3);
        \draw [hlddarrow] (s3) -- (d1);
        \draw [hlddarrow] (d1) -- node[above] {Yes} (a1);
        \draw [hlddarrow] (d1) -- node[left] {No} (s4);
        \draw [hlddarrow] (a1) |- (s1);

    \end{tikzpicture}
    \caption{System Functional Workflow}
\end{figure}

\subsection{System Failure Recovery Protocol}

Describe how the system detects, records, communicates, and recovers from
hardware, software, communication, or operational failures.
